\begin{minted}[bgcolor = lightgray, tabsize = 4, breaklines, numbers = left]{c}
    #include <stdio.h>
    #include <unistd.h>
    #include <fcntl.h>
    #include <string.h>
    #include <pthread.h>
    #include <arpa/inet.h>
    #include <stdlib.h>
    #include <stdint.h>
    #include <time.h>
    #include <math.h>
    #include <errno.h>
    #include <sys/socket.h>
    
    #include "Ring_Buffer.h"
    #include "sensor_read.h"
    
    #define PORT         50012
    #define CMD_BUF_SIZE 1024
    #define PI           3.14159265358979323846
    
    int simulation = 0;
    
    /* ============================================================================
     * GLOBAL SHUTDOWN FLAG + SOCKET FDS
     * Used by request_shutdown() to wake up all blocked calls and
     * signal every thread to exit cleanly.
     * ========================================================================= */
    static volatile int  g_shutdown      = 0;  /* set to 1 to begin shutdown     */
    static int           g_server_socket = -1; /* listening socket fd            */
    static int           g_client_socket = -1; /* currently connected client fd  */
    static pthread_mutex_t g_socket_lock = PTHREAD_MUTEX_INITIALIZER;
    
    /*
     * Called from any thread to trigger a full program shutdown.
     * Closes both sockets so that blocked accept() / recv() / send()
     * return immediately, allowing both threads to notice g_shutdown == 1
     * and return.
     */
    static void request_shutdown(void)
    {
        g_shutdown = 1;  /* visible to both threads before sockets are closed */
    
        pthread_mutex_lock(&g_socket_lock);
    
        if (g_client_socket >= 0) {
            shutdown(g_client_socket, SHUT_RDWR);
            close(g_client_socket);
            g_client_socket = -1;
        }
        if (g_server_socket >= 0) {
            shutdown(g_server_socket, SHUT_RDWR);
            close(g_server_socket);
            g_server_socket = -1;
        }
    
        pthread_mutex_unlock(&g_socket_lock);
    }
    
    /* ============================================================================
     * FAKE SENSOR DATA (simulation mode)
     * ========================================================================= */
    typedef struct {
        double   temp_base;
        double   adc0_base;
        double   adc1_base;
        uint64_t startup_time;
    } FakeSensorState;
    
    static FakeSensorState fake_state = {
        .temp_base    = 25,
        .adc0_base    = 325,
        .adc1_base    = 545,
        .startup_time = 0
    };
    
    double GenerateFakeTemperature(uint64_t elapsed_us)
    {
        (void)elapsed_us;
        return fake_state.temp_base;
    }
    
    uint32_t GenerateFakeADC0(uint64_t elapsed_us)
    {
        double t = elapsed_us / 1e6;
        return (uint32_t)(fake_state.adc0_base + 100.0 * sin(2.0 * PI * 3 * t));
    }
    
    uint32_t GenerateFakeADC1(uint64_t elapsed_us)
    {
        double t = elapsed_us / 1e6;
        return (uint32_t)(fake_state.adc1_base + 150.0 * cos(2.0 * PI * 5 * t));
    }
    
    uint8_t GenerateFakeSwitch(uint64_t elapsed_us)
    {
        static uint64_t last_change = 0;
        static uint8_t  state = 0x0F;
        if (elapsed_us - last_change > 5000000ULL) {
            state = (state == 0x0F) ? 0x00 : 0x0F;
            last_change = elapsed_us;
        }
        return state;
    }
    
    uint8_t GenerateFakePushButton(void)
    {
        return 100;
    }
    
    /* ============================================================================
     * THREAD STRUCTURES & GLOBALS
     * ========================================================================= */
    typedef struct {
        unsigned int temp_config;
        unsigned int adc0_config;
        unsigned int adc1_config;
        unsigned int switch_config;
        unsigned int pushb_config;
    } config_struct;
    
    typedef struct {
        pthread_t       thobj;
        pthread_cond_t  thcond;
        pthread_mutex_t thmutex;
        config_struct   th_config;
        short int       thkill;
    } thread_struct;
    
    typedef enum {
        sid_temp     = 1,
        sid_adc0     = 2,
        sid_adc1     = 3,
        sid_switches = 4,
        sid_pushb    = 5
    } sensor_id;
    
    typedef struct {
        int      sid;
        int      svalue;
        uint64_t ts;
    } packet;
    
    typedef enum {
        STOPPED       = 0,
        RUNNING,
        PAUSED,
        CONFIGURATION,
        DISCONNECT,
        SHUTDOWN          
    } stream_state_t;
    
    volatile stream_state_t stream_state;
    
    thread_struct sensor_thread;
    thread_struct server_thread;
    
    config_struct main_config = {
        .temp_config   = 300,
        .adc0_config   = 300,
        .adc1_config   = 300,
        .switch_config = 300,
        .pushb_config  = 300
    };
    
    /* ============================================================================
     * SENSOR THREAD
     * ========================================================================= */
    void *sensor_thread_function(void *arg)
    {
        (void)arg;
        int TempInitializerFlag = 0;
        printf("Sensor thread started...\n");
    
        int fd = open("/dev/meascdd", O_RDWR);
        if (fd < 0) {
            simulation = 1;
            printf("Device open failed — simulation mode\n");
        } else {
            printf("Device opened successfully\n");
        }
    
        InitTimer();
        InitRingBuffer();
    
        if (TempInit(fd) != 1) {
            printf("Temperature sensor init failed\n");
        } else {
            printf("Temperature sensor initialized\n");
            TempInitializerFlag = 1;
        }
    
        unsigned int adc0_data, adc1_data, dummy0, dummy1;
    
        uint64_t next_temp_time   = GetElapsedTime();
        uint64_t next_adc0_time   = GetElapsedTime();
        uint64_t next_adc1_time   = GetElapsedTime();
        uint64_t next_switch_time = GetElapsedTime();
        uint64_t next_pushb_time  = GetElapsedTime();
    
        uint64_t period_temp_us   = 1000000ULL / main_config.temp_config;
        uint64_t period_adc0_us   = 1000000ULL / main_config.adc0_config;
        uint64_t period_adc1_us   = 1000000ULL / main_config.adc1_config;
        uint64_t period_switch_us = 1000000ULL / main_config.switch_config;
        uint64_t period_pushb_us  = 1000000ULL / main_config.pushb_config;
    
        double t_adc0    = 0.0,  t_adc1    = 0.0;
        double freq_adc0 = 5.0,  freq_adc1 = 3.0;
    
        printf("Periods (us): T=%llu A0=%llu A1=%llu SW=%llu PB=%llu\n",
               (unsigned long long)period_temp_us,
               (unsigned long long)period_adc0_us,
               (unsigned long long)period_adc1_us,
               (unsigned long long)period_switch_us,
               (unsigned long long)period_pushb_us);
    
        while (1)
        {
            /* Check both thkill AND global shutdown */
            pthread_mutex_lock(&sensor_thread.thmutex);
            int kill = sensor_thread.thkill;
            pthread_mutex_unlock(&sensor_thread.thmutex);
            if (kill || g_shutdown) break;
    
            uint64_t current_time = GetElapsedTime();
    
            uint64_t sleep_target = next_temp_time;
            if (next_adc0_time   < sleep_target) sleep_target = next_adc0_time;
            if (next_adc1_time   < sleep_target) sleep_target = next_adc1_time;
            if (next_switch_time < sleep_target) sleep_target = next_switch_time;
            if (next_pushb_time  < sleep_target) sleep_target = next_pushb_time;
    
            if (current_time < sleep_target) {
                uint64_t sleep_us = sleep_target - current_time;
                while (sleep_us > 0) {
                    uint64_t chunk = (sleep_us > 100) ? 100 : sleep_us;
                    usleep(chunk);
                    sleep_us    -= chunk;
                    current_time = GetElapsedTime();
                    if (current_time >= sleep_target) break;
                    /* Exit sleep early if shutdown requested */
                    if (g_shutdown) goto sensor_exit;
                }
            }
    
            if (current_time >= next_temp_time) {
                uint32_t temp_data;
                if (simulation) {
                    temp_data = (uint32_t)(GenerateFakeTemperature(current_time) * 100);
                } else {
                    temp_data = TempRead3(fd);
                }
                PushRB(sid_temp, current_time, temp_data);
                next_temp_time += period_temp_us;
                printf("T=%u ", temp_data);
            }
    
            if (current_time >= next_adc0_time) {
                if (simulation) {
                    t_adc0   += 1.0 / main_config.adc0_config;
                    adc0_data = (uint32_t)(fake_state.adc0_base +
                                100.0 * sin(2.0 * PI * freq_adc0 * t_adc0));
                } else {
                    GetADC(fd, &adc0_data, &dummy0);
                }
                PushRB(sid_adc0, current_time, adc0_data);
                next_adc0_time += period_adc0_us;
                printf("A0=%u ", adc0_data);
            }
    
            if (current_time >= next_adc1_time) {
                if (simulation) {
                    t_adc1   += 1.0 / main_config.adc1_config;
                    adc1_data = (uint32_t)(fake_state.adc1_base +
                                150.0 * cos(2.0 * PI * freq_adc1 * t_adc1));
                } else {
                    GetADC(fd, &dummy1, &adc1_data);
                }
                PushRB(sid_adc1, current_time, adc1_data);
                next_adc1_time += period_adc1_us;
                printf("A1=%u ", adc1_data);
            }
    
            if (current_time >= next_switch_time) {
                uint8_t sw = simulation ? GenerateFakeSwitch(current_time)
                                        : ReadSwitch(fd);
                PushRB(sid_switches, current_time, sw);
                next_switch_time += period_switch_us;
                printf("SW=%u ", sw);
            }
    
            if (current_time >= next_pushb_time) {
                uint8_t pb = simulation ? GenerateFakePushButton()
                                        : ReadButton(fd);
                PushRB(sid_pushb, current_time, pb);
                next_pushb_time += period_pushb_us;
                printf("PB=%u ", pb);
            }
            printf("\n");
        }
    
    sensor_exit:
        if (fd >= 0) close(fd);
        (void)TempInitializerFlag;
        printf("Sensor thread exiting\n");
        return NULL;
    }
    
    /* ============================================================================
     * HELPERS
     * ========================================================================= */
    void set_config(config_struct *config)
    {
        pthread_mutex_lock(&server_thread.thmutex);
        main_config.temp_config   = config->temp_config;
        main_config.adc0_config   = config->adc0_config;
        main_config.adc1_config   = config->adc1_config;
        main_config.switch_config = config->switch_config;
        main_config.pushb_config  = config->pushb_config;
        pthread_mutex_unlock(&server_thread.thmutex);
    }
    
    void print_config(void)
    {
        printf("temp_config:     %u\n", main_config.temp_config);
        printf("adc0_config:     %u\n", main_config.adc0_config);
        printf("adc1_config:     %u\n", main_config.adc1_config);
        printf("switches_config: %u\n", main_config.switch_config);
        printf("pushb_config:    %u\n", main_config.pushb_config);
    }
    
    /* ============================================================================
     * SERVER THREAD
     * ========================================================================= */
    void *server_thread_function(void *arg)
    {
        (void)arg;
        packet payload[96];
        char   receive_command[CMD_BUF_SIZE];
    
        struct sockaddr_in server_address, client_addr;
        socklen_t client_addr_len = sizeof(client_addr);
    
        stream_state = PAUSED;
    
        /* Create listening socket */
        int server_sock = socket(AF_INET, SOCK_STREAM, 0);
        if (server_sock < 0) { perror("socket"); return NULL; }
    
        int opt = 1;
        setsockopt(server_sock, SOL_SOCKET, SO_REUSEADDR, &opt, sizeof(opt));
    
        memset(&server_address, 0, sizeof(server_address));
        server_address.sin_family      = AF_INET;
        server_address.sin_addr.s_addr = INADDR_ANY;
        server_address.sin_port        = htons(PORT);
    
        if (bind(server_sock, (struct sockaddr *)&server_address,
                 sizeof(server_address)) < 0) {
            perror("bind"); close(server_sock); return NULL;
        }
        if (listen(server_sock, 1) < 0) {
            perror("listen"); close(server_sock); return NULL;
        }
    
        /* Publish fd so request_shutdown() can close it */
        pthread_mutex_lock(&g_socket_lock);
        g_server_socket = server_sock;
        pthread_mutex_unlock(&g_socket_lock);
    
        printf("Server listening on port %d\n", PORT);
    
        /* ---- OUTER LOOP: wait for a client ---- */
        while (!g_shutdown)
        {
            printf("Waiting for client...\n");
    
            int client_sock = accept(server_sock,
                                     (struct sockaddr *)&client_addr,
                                     &client_addr_len);
            if (client_sock < 0) {
                if (g_shutdown) break;   /* accept() unblocked by request_shutdown() */
                perror("accept");
                continue;
            }
    
            /* Publish client fd */
            pthread_mutex_lock(&g_socket_lock);
            g_client_socket = client_sock;
            pthread_mutex_unlock(&g_socket_lock);
    
            printf("Client connected\n");
            stream_state = PAUSED;
    
            /* ---- INNER LOOP: serve one client ---- */
            while (!g_shutdown)
            {
                memset(receive_command, 0, CMD_BUF_SIZE);
                int n = recv(client_sock,
                             receive_command,
                             CMD_BUF_SIZE - 1,
                             MSG_DONTWAIT);
    
                if (n > 0)
                {
                    receive_command[n] = '\0';
                    for (char *p = receive_command; *p; p++)
                        if (*p == '\n' || *p == '\r') { *p = '\0'; break; }
    
                    printf("Command: [%s]\n", receive_command);
    
                    /* ---- start: one-shot batch send ---- */
                    if (!strcmp(receive_command, "start"))
                    {
                        int available;
                        pthread_mutex_lock(&server_thread.thmutex);
                        available = Ring_Buf.rb_cnt;
                        pthread_mutex_unlock(&server_thread.thmutex);
    
                        if (available > 0) {
                            int to_send = (available >= 96) ? 96 : available;
                            for (int i = 0; i < to_send; i++) {
                                if (PopRB(&payload[i].sid,
                                          &payload[i].ts,
                                          &payload[i].svalue) < 0) {
                                    to_send = i;
                                    break;
                                }
                            }
                            if (to_send > 0) {
                                ssize_t sent = send(client_sock, payload,
                                                    to_send * sizeof(packet), 0);
                                if (sent <= 0) { perror("send"); break; }
                                printf("Sent %d packets\n", to_send);
                            }
                        } else {
                            printf("start: ring buffer empty\n");
                        }
                        stream_state = PAUSED;
                    }
    
                    /* ---- pause / stop ---- */
                    else if (!strcmp(receive_command, "pause") ||
                             !strcmp(receive_command, "stop"))
                    {
                        stream_state = PAUSED;
                        printf("%s acknowledged\n", receive_command);
                    }
    
                    /* ---- configure ---- */
                    else if (!strncmp(receive_command, "configure ", 10))
                    {
                        config_struct cfg;
                        int parsed = sscanf(&receive_command[10],
                                            "%u %u %u %u %u",
                                            &cfg.temp_config,
                                            &cfg.adc0_config,
                                            &cfg.adc1_config,
                                            &cfg.switch_config,
                                            &cfg.pushb_config);
                        if (parsed == 5) { set_config(&cfg); print_config(); }
                        else printf("Bad configure arguments\n");
                    }
    
                    /* ---- disconnect: close this client, keep server alive ---- */
                    else if (!strcmp(receive_command, "disconnect"))
                    {
                        printf("Client requested disconnect\n");
                        stream_state = DISCONNECT;
                    }
    
                    /* ---- shutdown: tear down the entire server process ---- */
                    else if (!strcmp(receive_command, "shutdown"))
                    {
                        printf("Shutdown command received — stopping server\n");
                        stream_state = SHUTDOWN;
                        /*
                         * 1. Close client socket gracefully so the client
                         *    knows the connection ended cleanly.
                         * 2. Call request_shutdown() which sets g_shutdown = 1
                         *    and closes both sockets, unblocking accept() and
                         *    waking the sensor thread's sleep loop.
                         */
                        shutdown(client_sock, SHUT_RDWR);
                        close(client_sock);
    
                        pthread_mutex_lock(&g_socket_lock);
                        g_client_socket = -1;
                        pthread_mutex_unlock(&g_socket_lock);
    
                        request_shutdown();
    
                        /* Signal sensor thread to stop */
                        pthread_mutex_lock(&sensor_thread.thmutex);
                        sensor_thread.thkill = 1;
                        pthread_mutex_unlock(&sensor_thread.thmutex);
    
                        goto server_exit;  /* break out of both loops */
                    }
    
                    else
                    {
                        printf("Unknown command: [%s]\n", receive_command);
                    }
                }
                else if (n == 0)
                {
                    /* Peer closed connection normally */
                    printf("Client disconnected\n");
                    break;
                }
                else
                {
                    if (errno != EAGAIN && errno != EWOULDBLOCK && errno != EINTR) {
                        perror("recv");
                        break;
                    }
                }
    
                /* Handle disconnect state */
                if (stream_state == DISCONNECT) {
                    printf("Closing client connection\n");
                    shutdown(client_sock, SHUT_RDWR);
                    close(client_sock);
    
                    pthread_mutex_lock(&g_socket_lock);
                    g_client_socket = -1;
                    pthread_mutex_unlock(&g_socket_lock);
                    break;  /* go back to accept() */
                }
    
                usleep(1000);  /* light pacing */
            } /* inner loop */
    
            /* Clean up client socket if not already closed */
            pthread_mutex_lock(&g_socket_lock);
            if (g_client_socket >= 0) {
                close(g_client_socket);
                g_client_socket = -1;
            }
            pthread_mutex_unlock(&g_socket_lock);
    
        } /* outer loop */
    
    server_exit:
        /* Close listening socket if not already closed by request_shutdown() */
        pthread_mutex_lock(&g_socket_lock);
        if (g_server_socket >= 0) {
            close(g_server_socket);
            g_server_socket = -1;
        }
        pthread_mutex_unlock(&g_socket_lock);
    
        printf("Server thread exiting\n");
        return NULL;
    }
    
    /* ============================================================================
     * MAIN
     * ========================================================================= */
    int main(void)
    {
        printf("mng main running...\n");
    
        pthread_mutex_init(&server_thread.thmutex, NULL);
        pthread_mutex_init(&sensor_thread.thmutex, NULL);
        sensor_thread.thkill = 0;
    
        print_config();
    
        int rc;
    
        rc = pthread_create(&sensor_thread.thobj, NULL,
                            sensor_thread_function, NULL);
        if (rc != 0) {
            fprintf(stderr, "sensor thread failed: %s\n", strerror(rc));
            return 1;
        }
        printf("Sensor thread created\n");
    
        rc = pthread_create(&server_thread.thobj, NULL,
                            server_thread_function, NULL);
        if (rc != 0) {
            fprintf(stderr, "server thread failed: %s\n", strerror(rc));
            return 1;
        }
        printf("Server thread created\n");
    
        pthread_join(server_thread.thobj, NULL);
        pthread_join(sensor_thread.thobj, NULL);
    
        pthread_mutex_destroy(&sensor_thread.thmutex);
        pthread_mutex_destroy(&server_thread.thmutex);
        pthread_mutex_destroy(&g_socket_lock);
    
        printf("Program terminated cleanly\n");
        return 0;
    }
    \end{minted}